\SourcePage{433}{436}
\section{Transformation groups}

The collection of all symmetry transformations of a given body is called its 
group of symmetry transformations, or, more briefly, its \emph{symmetry 
group}.  Above, these transformations were described as geometrical motions 
of the body.  In quantum-mechanical applications, however, it is more 
convenient to regard them as transformations of the coordinates under which 
the Hamiltonian of the system remains invariant.  Clearly, if a rotation or 
reflection brings the system into coincidence with itself, the corresponding 
coordinate transformation leaves its Schr\"odinger equation unchanged.  We 
shall therefore speak of the group of transformations under which the given 
Schr\"odinger equation is invariant\footnote{This viewpoint permits us to 
include not only the rotation and reflection groups considered here, but also 
other kinds of transformations that leave the Schr\"odinger equation 
unchanged.  Among them are permutations of the coordinates of identical 
particles belonging to the system (a molecule or an atom).  The collection of 
all permutations of identical particles that are possible in a given system 
is called its \emph{permutation group}; such groups have already occurred in 
\S~63.  The general group properties developed below apply to permutation 
groups as well, but we shall not pursue a more detailed study of groups of 
this kind.

One comment is needed concerning the notation used in this chapter.  Symmetry 
transformations are, in substance, operators of the same kind as those used 
throughout the book, and so their letters ought to carry hats.  We omit the 
hats in deference to customary notation and because no ambiguity can result 
here.  For the same reason the identity transformation is denoted by the 
usual symbol $E$, rather than by 1 as the notation of the other chapters would 
suggest.  Finally, in this chapter the inversion operator is written $I$, in 
place of the symbol $P$ used in \S~30 and customary in present-day quantum-
mechanics literature.}.


Symmetry groups are conveniently studied by means of the general mathematical machinery known as \emph{group theory}, whose foundations are presented below. We begin with groups containing only a finite number of distinct transformations (the so-called \emph{finite groups}). Each transformation belonging to a group is called an \emph{element of the group}.

Symmetry groups have the following evident properties. Every group contains the identity transformation $E$ (called the \emph{identity element of the group}). Group elements can be \emph{multiplied} together; the product of two (or several) transformations means the result of applying them successively. Evidently, the product of any two elements of a group is an element of the same group. Element multiplication obeys the associative law $(AB)C=A(BC)$, where $A,B,C$ are group elements. The commutative law does not hold in general, however: ordinarily $AB\ne BA$. For every group element $A$, the same group contains an \emph{inverse element} $A^{-1}$ (the inverse transformation), such that $AA^{-1}=E$. In some
\SourcePage{435}{438}
cases an element may coincide with its inverse; in particular, $E^{-1}=E$. It is evident that mutually inverse elements $A$ and $A^{-1}$ commute.

The inverse of the product $AB$ of two elements is
\[
(AB)^{-1}=B^{-1}A^{-1}
\]
and the analogous statement holds for a product of more elements. This follows at once by carrying out the multiplication and using associativity.

If all the elements of a group commute, the group is called \emph{abelian}. The so-called \emph{cyclic groups} form a special class of abelian groups. A cyclic group is one whose every element can be obtained by taking successive powers of a single element; thus it consists of the elements
\[
A,A^2,A^3,\ldots,A^n=E,
\]
where $n$ is an integer.

Let $\boldsymbol{G}$ be a group\footnote{Groups will be denoted by bold italic letters.}. If a collection of its elements, denoted by $\boldsymbol{H}$, can be selected so that it too forms a group, then $\boldsymbol{H}$ is called a \emph{subgroup} of $\boldsymbol{G}$. One and the same group element may belong to several different subgroups.

Take any element $A$ of the group and form its successive powers. Since the total number of elements in the group is finite, the identity element must eventually be reached. If $n$ is the smallest number for which $A^n=E$, then $n$ is called the \emph{order of the element} $A$, while the collection $A,A^2,\ldots,A^n=E$ is the \emph{period} of $A$. The period is denoted by $\{A\}$; it is itself a group, hence a subgroup of the original group, and specifically a cyclic subgroup.

To determine whether a given collection of group elements is a subgroup, it is enough to check that multiplication of any two of its elements gives an element contained in the same collection. Indeed, together with every element $A$, the collection will then contain all its powers, including $A^{n-1}$ (where $n$ is the order of the element), which serves as its inverse (since $A^{n-1}A=A^n=E$); it will plainly contain the identity as well.

The total number of elements in a group is called the \emph{order} of the 
group.  It is readily seen that the order of a subgroup divides the order of 
the whole
\SourcePage{436}{439}
group.  To prove this, consider a subgroup $\boldsymbol{H}$ of a group 
$\boldsymbol{G}$, and let $G_1$ be an element of $\boldsymbol{G}$ that does not 
belong to $\boldsymbol{H}$.  Multiplication of every element of 
$\boldsymbol{H}$ by $G_1$---on the right, for example---gives a set (or 
\emph{coset}) of elements denoted by $\boldsymbol{H}G_1$.  Every element of 
this coset plainly belongs to $\boldsymbol{G}$, but none belongs to 
$\boldsymbol{H}$.  Indeed, if two elements $H_a,H_b$ of $\boldsymbol{H}$ 
satisfied $H_aG_1=H_b$, it would follow that 
$G_1=H_a^{-1}H_b$; this would put $G_1$ in $\boldsymbol{H}$, contrary to our 
assumption.  In the same way, if $G_2$ belongs to $\boldsymbol{G}$ but to 
neither $\boldsymbol{H}$ nor $\boldsymbol{H}G_1$, then none of the elements of 
$\boldsymbol{H}G_2$ belongs to either of those two sets.  Continuing in this 
fashion eventually exhausts all elements of the finite group 
$\boldsymbol{G}$.  Thus its elements are partitioned into sets (the 
\emph{cosets} of $\boldsymbol{H}$ in $\boldsymbol{G}$)
\[
\boldsymbol{H},\ \boldsymbol{H}G_1,\ \boldsymbol{H}G_2,\ldots,\boldsymbol{H}G_m,
\]
each containing $h$ elements, where $h$ is the order of the subgroup 
$\boldsymbol{H}$.  Consequently the order of $\boldsymbol{G}$ is $g=hm$, 
which proves the assertion.  The integer $m=g/h$ is called the \emph{index} 
of $\boldsymbol{H}$ in $\boldsymbol{G}$.


If the order of a group is prime, the result just proved immediately shows that the group has no subgroups at all (apart from $E$ and the group itself). The converse is also true: any group without subgroups must have prime order and must moreover be cyclic (otherwise it would contain elements whose periods would form subgroups).

We now introduce the important notion of \emph{conjugate elements}.  Two 
elements $A$ and $B$ are said to be conjugate to one another if
\[
A=CBC^{-1},
\]
where $C$ is also an element of the group.  Multiplying this equality by $C$ 
on the right and by $C^{-1}$ on the left gives the inverse relation 
$B=C^{-1}AC$.  An essential property of conjugacy is transitivity: if $A$ is 
conjugate to $B$, and $B$ to $C$, then $A$ and $C$ are conjugate as well.  In 
fact, from $B=P^{-1}AP$ and $C=Q^{-1}BQ$, with $P,Q$ elements of the group, it 
follows that $C=(PQ)^{-1}A(PQ)$.  One may therefore consider sets of elements 
that are mutually conjugate.  Such sets are called \emph{conjugacy classes}, 
or simply classes of the group.  Any one of its elements $A$ determines the 
entire class: once $A$ is given, form the products $GAG^{-1}$ as $G$ ranges 
over all elements of the group; the same member of the class may, of course, 
be obtained more than once.  In this way the whole group can be partitioned 
into classes, and each element evidently belongs to only one class.
\SourcePage{437}{440}
The identity element forms a class by itself, since $GEG^{-1}=E$ for every 
element $G$ of the group.  In an Abelian group the same is true of each 
element.  All elements of such a group commute by definition, so every 
element is conjugate only to itself and hence constitutes its own class.  It 
should be emphasized that a class other than $E$ is by no means a subgroup: 
this already follows from the fact that it does not contain the identity 
element.


All elements in the same class have the same order. Indeed, if $n$ is the order of $A$ (so that $A^n=E$), then for the conjugate element $B=CAC^{-1}$ we likewise have $(CAC^{-1})^n=CA^nC^{-1}=E$.

Let $\boldsymbol{H}$ be a subgroup of $\boldsymbol{G}$, and let $G_1$ be an element of $\boldsymbol{G}$ outside $\boldsymbol{H}$. It is easily verified that the collection of elements $G_1\boldsymbol{H}G_1^{-1}$ satisfies all the requirements for a group and hence is also a subgroup of $\boldsymbol{G}$. The subgroups $\boldsymbol{H}$ and $G_1\boldsymbol{H}G_1^{-1}$ are called conjugate; each element of one is conjugate to an element of the other. Giving $G_1$ different values yields a number of conjugate subgroups, some of which may coincide. It can happen that every subgroup conjugate to $\boldsymbol{H}$ coincides with $\boldsymbol{H}$. In that case $\boldsymbol{H}$ is called a \emph{normal divisor} (or an \emph{invariant subgroup}) of $\boldsymbol{G}$. Thus, for example, every subgroup of an abelian group is plainly a normal divisor of that group.

Consider a group $\boldsymbol{A}$ with $n$ elements $A,A',A'',\ldots$ and a group $\boldsymbol{B}$ with $m$ elements $B,B',B'',\ldots$, and suppose that all the elements of $\boldsymbol{A}$ (apart from the identity $E$) are distinct from and commute with the elements of $\boldsymbol{B}$. Multiplying each element of $\boldsymbol{A}$ by every element of $\boldsymbol{B}$ produces a collection of $nm$ elements that likewise forms a group. Indeed, for any two elements of the collection,
$AB\cdot A'B'=AA'\cdot BB'=A''B''$, which is again an element of the same collection. The resulting group of order $nm$ is denoted by $\boldsymbol{A}\times\boldsymbol{B}$ and is called the \emph{direct product} of the groups $\boldsymbol{A}$ and $\boldsymbol{B}$.

Finally, we introduce the concept of group \emph{isomorphism}. Two groups $\boldsymbol{A}$ and $\boldsymbol{B}$ of the same order are called isomorphic if a one-to-one correspondence can be established between their elements such that, if $A$ corresponds to $B$
\SourcePage{438}{441}
and $A'$ to $B'$, then $A''=AA'$ corresponds to $B''=BB'$. Considered abstractly, two such groups plainly have identical properties, even though the concrete meanings of their elements differ.
